\section{Conclusion and Discussion}
\label{sec:conclusion}

In this work, we introduced \textbf{QP-Merge}, a training-free, quantization-preserving model merging framework designed specifically to address the severe performance degradation that occurs when merging task-specific foundation models for low-bit edge deployment. Grounded in a pragmatic research philosophy, QP-Merge reconciles the conflict between floating-point parameter merging and hardware-accelerated integer execution. By decomposing task vectors into a dense quantized path and an ultra-sparse high-precision outlier path (Outlier-Residual Decoupling), we successfully insulate symmetric uniform quantization grids from range stretching. Furthermore, by introducing a rapid, zero-labeled-data activation reconstruction step (Quantization-Error Aware Scale Calibration), we correct multi-task activation scale mismatches directly in the quantized regime.

Our extensive empirical evaluations on dual-task vision classification (MNIST and SVHN) utilizing a pre-trained \texttt{ViT-B-32} base model prove that QP-Merge completely resolves the low-bit quantization bottleneck. In 4-bit (INT4) mode, QP-Merge achieves an average accuracy of \textbf{94.70\%} across 3 random seeds, recovering over \textbf{88\%} of the performance drop caused by naive quantization relative to the optimized baseline and performing within \textbf{0.42\%} of the optimized unquantized FP32 merged bound. In 8-bit (INT8) mode, our unsupervised calibration aligns activation scales and simultaneously optimizes task coefficients on-the-fly, achieving an average accuracy of \textbf{95.14\%} which represents a \textbf{+0.02\%} gain over the optimized unquantized FP32 baseline (and is virtually lossless compared to unquantized references).

\subsection{Limitations and Future Work}
While our framework exhibits strong robustness and hardware-friendly characteristics, there are certain areas we identify for future exploration:
\begin{itemize}
    \item \textbf{Outlier Density and Multi-Task Scaling:} Standard multi-task merging benchmarks evaluate on 8-task image classification suites or multi-task NLP benchmarks. When scaling QP-Merge from $T=2$ tasks to $T=8$ or more, two key scaling phenomena arise:
    \begin{enumerate}
        \item \emph{Outlier Overlap and Density Control:} In ORD, we isolate the top $1\%$ (or $0.5\%$) of extreme weight differences as outliers. If task-specific outliers are completely disjoint, the union of outliers across $T=8$ tasks could theoretically scale up to $8\%$, which would degrade sparse matrix multiplication (SpMM) efficiency and increase VRAM footprint. However, in practice, weight outliers in neural network layers often cluster around highly critical feature-projection pathways (e.g., specific channels inside key, query, or projection matrices). This leads to substantial spatial overlap in outlier locations across different tasks. To manage outlier density rigorously at scale, we propose employing a \emph{global thresholding scheme} (enforcing a global percentile constraint on the merged weight matrix $\sum \lambda_t \Delta W_t$) rather than individual per-task extraction. This ensures that the combined outlier density is strictly bounded by a user-defined threshold (e.g., $1.0\%$), regardless of the number of merged tasks.
        \item \emph{Compounding Activation Scale Mismatch:} As more tasks are merged, the internal scale distributions become increasingly complex. Joint optimization of layer-wise diagonal scalers $D_l$ and blending parameters $\lambda$ via QE-Calib over $M=128$ domain-balanced calibration samples remains highly data-efficient. The computational complexity of QE-Calib scales linearly with the number of layers but remains constant with respect to the number of tasks, as it is a zero-labeled-data reconstruction process on final embeddings. Enforcing domain balancing in the calibration dataset (e.g., sampling $128 / T$ images per task) ensures that the optimization treats each task distribution equitably, preventing scale drift on minority domains.
    \end{enumerate}
    \item \textbf{Edge Hardware Heterogeneous Compatibility:} Many low-power IoT microcontrollers and legacy edge NPUs only support homogeneous fixed-point execution (pure INT4 or INT8) and lack floating-point execution units or the ability to run dynamic, hybrid representations. To support strictly integer-only platforms, QP-Merge can be extended to quantize outliers to a slightly higher bit-width (such as INT8 or INT16) while keeping the dense remainder in INT4, or by pruning the outliers entirely for strict homogeneous INT4 execution.
    \item \textbf{Sparse Runtime Libraries:} Although SpMM is highly efficient on modern sparse tensor cores (e.g., NVIDIA Ampere/Ada Lovelace architectures), actual speedups depend on compiler-level support and standard library integration (such as CUTLASS or Triton). Standardizing sparse layouts across diverse edge chips is an ongoing hardware effort.
    \item \textbf{Scaling to Autoregressive Large Language Models (LLMs):} While our evaluation focused on a vision transformer block, LLMs are known to exhibit even more severe activation outlier channels \cite{dettmers2022gpt3}. Extending QP-Merge to billions-of-parameters autoregressive language networks remains a high-impact future direction.
\end{itemize}

\subsection{Practical Broader Impact}
From a deployment perspective, QP-Merge drastically reduces VRAM bandwidth and serving costs (achieving a \textbf{3.77$\times$ weight compression ratio} in INT4) while retaining the zero-cost advantage of model merging. By enabling near-lossless multi-task capabilities on commodity edge devices under strict memory and latency SLAs without requiring labeled training datasets, QP-Merge lowers the barrier for deploying powerful foundation models in real-world, localized consumer applications.
